Our research relates to the broad literature on the political effect of immigration. Much of the existing literature emphasizes recent immigration trends and their immediate impact on public opinion and political behavior \citep{halla_immigration_2017, dustmann_refugee_2018, steinhardt_xenophobic_2018}. \citet{steinmayr_contact_2021} demonstrates that the type of contact matters: in Upper Austria, temporary refugee passage towards Germany in 2015 increased far-right voting, whereas longer-term settlement reduced it. In the French context, studies have found that larger immigrant populations tend to correlate with stronger electoral performance for the Front National at the département level \citep{edo_immigration_2019}. However, this pattern does not hold at the more localized municipal level, where the relationship appears to be reversed \citep{della2013competitive, vasilopoulos2022residential, vertier_dismantling_2023}. \citet{dazey_mosque_2024} finds that support for the Front National tends to rise in polling stations located at moderate distances from mosques, but declines beyond that range. 

A key feature of the natural experiment we draw on is that the cultural distance between the forced immigrants and the local population was much smaller than in other contexts. Meanwhile, the refugee flow was highly politically selected, which allows for an exploration of its effects on political attitudes not only through backlash or contact, but also through the diffusion of ideas. In this way, our study contributes to the literature on migration and the transmission of political ideas \citep{spilimbergo_democracy_2009, barsbai_effect_2017, dippel_leadership_2021, calderon_racial_2023, bazzi_confederate_2023}. This literature concentrates mostly on the effect at the origin country, with the exception of \citet{dippel_leadership_2021}, who study how the leaders of the failed German revolution of 1848–1849, expelled to the United States, positively influenced antislavery culture. Since our research is framed within a democratic context, we can distinguish ourselves by analysing the effect on electoral outcomes and controlling for previous political affiliations. 

%Several empirical studies provide support for the “contact hypothesis” (Allport, 1954) by showing that contact with minorities can reduce cultural anxiety and prejudice, thereby mitigating racial bias (Schindler&Westcott, 2021) and aecting partisanship (Billings et al., 2021). Similarly, sustained interaction with immigrants and personal family history of immigration reduces hostility and far-right voting (Dinas et al., 2021; Steinmayr, 2021) and increases altruistic behavior toward the group to which individuals are exposed (Bursztyn et al., 2024). From Lang & Schneider, 2024. 